% Part: methods
% Chapter: proofs
% Section: starting-proofs

\documentclass[../../../include/open-logic-section]{subfiles}

\begin{document}

\olfileid{mth}{prf}{str}

\olsection{د ثبوت پيل}

خو له کومه ځايه پيل کوئ؟

تاسو ته د ثابتولو دپاره يوه خبره درکړل شوې؛ نو همدا
بايد د ثبوت په پای کښې وروستۍ خبره وي (البته په سر
کښې \emph{اعلانولے} شئ چې دا به ثابتوئ، خو د فرض
په توګه يې نه شئ کارولے)۔ هغه څه چې ثابتول يې
غواړئ د کاغذ د نوې پاڼې په ښکته برخه کښې وليکئ،
څو مو هدف له سترګو پناه نه شي۔

بيا ښايي ځينې داسې فرضونه ولرئ چې کارولے يې شئ
(دا خبره به په راتلونکې برخه کښې، د خپل ثبوت د
\emph{ډول} د بحث پر مهال، نوره هم روښانه شي)۔ دا
فرضونه د پاڼې په سر کښې وليکئ او څرګند کړئ چې
فرضونه دي (يعنې که $p$ فرضوئ، وليکئ «فرض کړئ
چې $p$»، يا «ونيسئ چې $p$»)۔ په پای کښې ښايي
په پوښتنه کښې داسې تعريفونه وي چې پېژندل يې پکار
دي۔ کېدے شي د يوه ځانګړي تعريف د کارولو غوښتنه
درڅخه شوې وي، يا په فرضونو يا هغه پايله کښې چې
ورته رسېدل غواړئ بېلابېل تعريفونه راغلي وي۔
\emph{دا وليکئ او ډاډ ترلاسه کړئ چې په معنا يې پوهېږئ۔}

د ثبوت د برابرولو بڼه د پوښتنې په شکل هم تړلې ده۔
راتلونکې برخه ښيي چې د جملې د ډول له مخې خپل
ثبوت څنګه برابر کړئ۔

\end{document}
